\documentclass[11pt]{article}
\usepackage[margin=0.5in]{geometry}
\usepackage{enumitem}
\usepackage{hyperref}
\usepackage{titlesec}
\titleformat{\section}{\large\bfseries}{}{0em}{}
\setlength{\parskip}{0.2em}
\setlength{\parindent}{0em}

\begin{document}

\begin{center}
    {\LARGE \textbf{Evgeny Redekop}} \\
    Thouless Postdoctoral Fellow, University of Washington \\
    Email: \href{mailto:eugene.redekop@gmail.com}{eugene.redekop@gmail.com} \quad | \quad +1 (805) 568-8958 \\
    \href{https://evgeny-redekop.github.io}{Website} \quad | \quad \href{https://scholar.google.com/citations?hl=en&user=eXBKws4AAAAJ}{Google Scholar} \quad | \quad \href{https://github.com/evgeny-redekop}{GitHub} \quad | \quad \href{https://www.linkedin.com/in/evgeny-redekop/}{LinkedIn}
\end{center}

\section*{Research Interests}
Quantum instrumentation bridging condensed matter and quantum information; nanoSQUID-on-tip and Josephson microwave microscopy; superconducting devices and gate-tunable Josephson-junction arrays; spatially resolved spectroscopy of correlated and topological quantum materials; cryogenic low-noise measurement, laboratory automation, and machine-learning--assisted data analysis.

\section*{Education}
\textbf{University of California, Santa Barbara} \hfill 2021--2026\\
Ph.D. in Physics\\
Advisor: Prof. Andrea F. Young\\
Dissertation: \emph{Local magnetic imaging of correlated states in moir\'e systems using nanoSQUID-on-tip microscopy}

\vspace{0.3em}
\textbf{Skolkovo Institute of Science and Technology} \hfill 2018--2020\\
M.Sc. in Materials Science

\vspace{0.3em}
\textbf{Lomonosov Moscow University} \hfill 2014--2018\\
B.Sc. in Physics

\section*{Research Experience}

\textbf{Thouless Postdoctoral Fellow} \hfill 2026 -- Present\\
\emph{University of Washington (groups of Charlie Marcus and Matthew Yankowitz)}
\begin{itemize}[leftmargin=*]
    \item Studying dissipation-tuned superconductor--insulator transitions and anomalous-metal physics in gate-defined Al/InAs Josephson-junction arrays at dilution-refrigerator temperatures.
    \item Building human-in-the-loop software, including LLM-based agents, for measurement planning and collaborative data analysis.
\end{itemize}

\vspace{0.3em}
\textbf{Graduate Researcher, NSF Quantum Foundry} \hfill Jan 2021 -- 2026\\
\emph{University of California, Santa Barbara (Advisor: Andrea F. Young)}
\begin{itemize}[leftmargin=*]
    \item Developed and optimized \textbf{nanoSQUID-on-tip (nSOT)} probes, improving magnetic flux sensitivity by over 50× to 4 n$\Phi_0/\sqrt{Hz}$, the highest reported to date (manuscript in preparation).
    \item Designed custom \textbf{UHV cryogenic deposition tools} and Python-based multi-instrument control software.
    \item Performed first \textbf{local magnetic imaging of fractional Chern insulators} in twisted MoTe$_2$ (\textit{Nature}, 2024).
    \item Discovered \textbf{current-induced spin--orbit torque switching} in moir\'e Chern magnets (\textit{Nature Physics}, 2023).
    \item Collaborated with theorists and experimentalists (L. Fu, A. MacDonald, L. Levitov, L. Glazman, K. Mak, J. Shan, X. Xu, S. Wilson, C. Palmstrøm).
    \item Supervised several undergraduate projects on cryogenic nSOT system development and fabrication yield optimization.
\end{itemize}

\textbf{Scanning Electron Microscopy Superuser} \hfill Oct 2022 -- May 2024\\
University of California, Santa Barbara\\
Trained and supported new users on Thermo Fisher Apreo C SEM system.

\textbf{Student Researcher, Skolkovo Institute of Science and Technology} \hfill Oct 2018 -- Oct 2020\\
Advisor: Prof. Albert G. Nasibulin\\
Built experimental setup for synthesis and characterization of metal--organic frameworks; conducted STM, FTIR, and THz spectroscopy measurements; developed data analysis pipelines in Python.

\section*{Awards and Fellowships}
\begin{itemize}[leftmargin=*]
    \item \textbf{Thouless Postdoctoral Fellowship} (2026), University of Washington -- Competitive named postdoctoral fellowship in physics.
    \item \textbf{NSF Quantum Foundry Fellowship} (2021--2025) -- Four-year competitive fellowship supporting quantum materials instrumentation.
    \item \textbf{Art of Science Competition -- Honorable Mention} (2023), UCSB.
    \item \textbf{FASIE Grant for Young Entrepreneurs} (2015) -- Foundation for Assistance to Small Innovative Enterprises.
    \item \textbf{Diploma with Distinction}, Skolkovo Institute of Science and Technology (2020).
\end{itemize}

\section*{Technical Skills}
Cryogenic scanning probe microscopy; characterization tools (AFM, SEM); UHV thin-film deposition (In, Pb); low-noise measurement system design; vacuum system assembly and debugging; nanofabrication and materials characterization; Python/MATLAB scientific computing for instrumentation control and analysis; LLM- and agent-based tooling for laboratory automation and data analysis.

\section*{Publications}
\begin{itemize}[leftmargin=*]
    \item C. Zhang*, \textbf{E. Redekop}*, H. Stoyanov, J. H. Farrell, S. Kim, L. Holleis, D. Gong, A. Keough, Y. Choi, T. Taniguchi, K. Watanabe, M. E. Huber, A. C. Bleszynski Jayich, A. Lucas, A. F. Young. \textit{Imaging flat band electron hydrodynamics in biased bilayer graphene}, arXiv:2603.11175 (2026)
    \item W. Li*, \textbf{E. Redekop}*, C. Wang Beach*, C. Zhang*, X. Zhang, X. Liu, W. Holtzmann, C. Hu, E. Anderson, H. Park, T. Taniguchi, K. Watanabe, J.-H. Chu, L. Fu, T. Cao, D. Xiao, A. Young, X. Xu. \textit{Universal Magnetic Phases in Twisted Bilayer MoTe$_2$}, Nano Letters 25, 18044--18050 (2025)
    \item Caitlin L. Patterson, Owen I. Sheekey, Trevor B. Arp, Ludwig F. W. Holleis, Jin Ming Koh, Youngjoon Choi, Tian Xie, Siyuan Xu, Yi Guo, Hari Stoyanov, \textbf{Evgeny Redekop}, Canxun Zhang, Grigory Babikyan, David Gong, Haoxin Zhou, Xiang Cheng, Takashi Taniguchi, Kenji Watanabe, Martin E. Huber, Chenhao Jin, Étienne Lantagne-Hurtubise, Jason Alicea, Andrea F. Young, \textit{Superconductivity and spin canting in spin–orbit-coupled trilayer graphene}, Nature, 641, 632–638 (2025)
    \item Ganesh Pokharel, Canxun Zhang, \textbf{Evgeny Redekop}, Brenden R Ortiz, Andrea N Salinas, Sarah Schwarz, Steven J Gomez Alvarado, Suchismita Sarker, Andrea F Young, Stephen D Wilson, \textit{Evolution of charge correlations in the hole-doped kagome superconductor CsV$_{3-x}$Ti$_x$Sb$_5$"}, Physics Review Materials (2025)
    \item \textbf{E. Redekop}, C. Zhang, H. Park, J. Cai, E. Anderson, O. Sheekey, T. Arp, G. Babikyan, S. Salters, K. Watanabe, T. Taniguchi, M. Huber, X. Xu, A. Young. \textit{Direct magnetic imaging of fractional Chern insulators in twisted MoTe$_2$}, Nature 635, 584–589 (2024).
    \item Rajarshi Bhattacharyya, \textbf{Evgeny Redekop}, Mirko Bacani, \textit{Phases of Chern insulator in moiré bilayers}, White paper, attocube systems AG (2024)
    \item Trevor Arp, Owen Sheekey, Haoxin Zhou, CL Tschirhart, Caitlin L Patterson, HM Yoo, Ludwig Holleis, \textbf{Evgeny Redekop}, Grigory Babikyan, Tian Xie, Jiewen Xiao, Yaar Vituri, Tobias Holder, Takashi Taniguchi, Kenji Watanabe, Martin E Huber, Erez Berg, Andrea F Young, \textit{Intervalley coherence and intrinsic spin-orbit coupling in rhombohedral trilayer graphene}, Nature Physics,  20, 1413–1420 (2024)
    \item Wenjin Zhao, Kaifei Kang, Lizhong Li, Charles Tschirhart, \textbf{Evgeny Redekop}, Kenji Watanabe, Takashi Taniguchi, Andrea Young, Jie Shan, Kin Fai Mak, \textit{Realization of the Haldane Chern insulator in a moiré lattice}, Nature Physics, 20, 275–280 (2024)
    \item C.L. Tschirhart*, \textbf{E. Redekop}*, L. Li, T. Li, S. Jiang, T. Arp, O. Sheekey, T. Taniguchi, K. Watanabe, M. Huber, K. Mak, J. Shan, A. Young. \textit{Intrinsic spin Hall torque in a moir\'e Chern magnet}, Nature Physics 19, 807–813 (2023).
    \item Victor I. Kleshch, Alexander A. Tonkikh, Sergey A. Malykhin, \textbf{Eugene V. Redekop}, Andrey S. Orekhov, Andrey L. Chuvilin, Elena D. Obraztsova, and Alexander N. Obraztsov, \textit{Field emission from single-walled carbon nanotubes modified by annealing and CuCl doping}, Applied Physics Letters 109, 143112 (2016)
\end{itemize}

\section*{Mentorship}
Supervised undergraduate researchers at UCSB on projects including nSOT fabrication optimization and cryothermal evaporator design (2021–2025). Alumni have continued to graduate programs at UCSB and UCSD.

\section*{Public Talks}
\begin{itemize}[leftmargin=*]
    \item APS March Meeting 2025 -- \textit{Local imaging of ferromagnetism and topology in twisted bilayer MoTe$_2$}
    \item APS March Meeting 2024 -- \textit{Magnetic imaging of integer and fractional Chern insulators in tMoTe$_2$}
    \item APS March Meeting 2023 -- \textit{Intrinsic spin Hall torque in a moir\'e Chern magnet}
\end{itemize}

\end{document}
